%======================================================================
\section{The counting criterion and sequential extraction}\label{app:imported}
%======================================================================

The lower bounds of Theorems~A and~B stand on two results of the companion papers: the counting criterion
(Theorem~\ref{thm:criterion}, which is~\cite[Theorem~8.1]{DouglasMghirbiA}) and sequential extraction (Theorem~\ref{thm:extraction},
which is~\cite[Theorem~5.2]{DouglasMghirbiB}). This appendix reproduces both proofs in the notation of this paper, so that a reader can check
the main theorems without the companions. The remaining imports, chiefly the Dehn bound for $H_3$ (Proposition~\ref{prop:dehn}), are listed in
Table~\ref{tab:dependencies}. Within this appendix, $\|uv-vu\|_2$ measures how far two unitaries are from commuting, and $\ln$ is the natural
logarithm.

\subsection{Three facts about copies of \texorpdfstring{$S_3$}{S3}}

The first fact rounds an approximate representation of $S_3=\langle s,t\mid s^2,t^2,(st)^3\rangle$ to an exact one.

\begin{lemma}[Exact $S_3$ rounding]\label{lem:imp-s3round}
If $A,B\in U(D)$ satisfy $\|A^2-I\|_2,\|B^2-I\|_2,\|(AB)^3-I\|_2\le e$, there is a unitary representation $u$ of $S_3$ on $\mathbb C^D$
with $\|u(s)-A\|_2\le e$, $\|u(t)-B\|_2\le8e$ and $\|u(st)-AB\|_2\le9e$, and each normal form $1,s,t,st,ts,sts$ evaluated in $u$ is within $10e$
of the same word in $A,B$.
\end{lemma}

\begin{proof}
Round the spectra of $A$ and $B$ to the nearer of $\pm1$. For $|z|=1$ this moves $z$ by at most $|z^2-1|$, so the resulting Hermitian unitaries
$S,T$ satisfy $\|S-A\|_2,\|T-B\|_2\le e$. Then $C=ST$ satisfies $SCS=C^*$ and $\|C^3-I\|_2\le7e$. Round the spectrum of $C$ to the nearest cube root
of unity, choosing $1$ at ties, with one exception: send $-1$ to $1$, although it is not the nearest root there. At that point the
displacement is $2=|(-1)^3-1|$. This rule commutes with complex conjugation and moves $z$ by at most $|z^3-1|$, so $C_0$ satisfies $C_0^3=I$,
$SC_0S=C_0^*$ and $\|C_0-C\|_2\le7e$. Then $u(s)=S$ and $u(t)=SC_0$ are involutions with $u(s)u(t)=C_0$, and the bounds follow from the triangle
inequality.
\end{proof}

The second fact is the entropy that nearly commuting exact copies force. Write $h_2$ for the binary entropy.

\begin{lemma}[Entropic sector bound]\label{lem:imp-sector}
Let $u_1,\dots,u_N:S_3\to U(D)$ be unitary representations, not assumed to commute, such that for a three-cycle $c$
\[
\|u_i(c)-I\|_2\ge\Delta\quad(1\le i\le N),\qquad\|u_i(g)u_j(h)-u_j(h)u_i(g)\|_2\le\eta\quad(i\ne j,\ g,h\in S_3).
\]
Put $\theta=\sqrt{2N}\,\eta$. If $\theta\le\Delta$ and $\theta^2\le5/6$, then
\[
\log D\ \ge\ N\Big[\frac{(\Delta-\theta)^2}3-h_2(\theta^2)-\theta^2\log5\Big].
\]
\end{lemma}

\begin{proof}
Let $\Sigma=S_3^N$ and $f(g)=u_1(g_1)\cdots u_N(g_N)$. For $h$ in factor $i$, $\|f(gh)-f(g)u_i(h)\|_2=\|Y^*XY-X\|_2$ with $X=u_i(h)$ and
$Y=u_{i+1}(g_{i+1})\cdots u_N(g_N)$, and $\|Y^*XY-X\|_2^2=2-2\,\mathrm{Re}\,\tau(X^*Y^*XY)$. The averages
$T_j(Z)=\frac16\sum_au_j(a)^*Zu_j(a)$ are orthogonal projections on $M_D$ with the normalized Hilbert--Schmidt inner product, the average of
$Y^*XY$ over $g$ is $T_N\cdots T_{i+1}(X)$, and $\|(I-T_j)X\|_2^2=\frac1{12}\sum_a\|[X,u_j(a)]\|_2^2\le\eta^2/2$.
Lemma~\ref{lem:projection-product} therefore gives
\[
\mathbb E_g\|f(gh)-f(g)u_i(h)\|_2^2\le2(N-i)\eta^2\le\theta^2 .
\]
Let $\mathsf R$ be the right regular representation of $\Sigma$ and $W:\mathbb C^D\to\ell^2(\Sigma)\otimes\mathbb C^D$ the isometry
$Wv=|\Sigma|^{-1/2}\sum_g|g\rangle\otimes f(g)v$. Then $\|\mathsf R(h)W-Wu_i(h)\|_{HS}\le\sqrt D\,\theta$ for $h$ in factor $i$. Put $\Pi=WW^*$ and
$\rho=\Pi/D$, a state of entropy $\log D$. For the three-cycle $c_i$ of factor $i$ and $Z_i=2I-\mathsf R(c_i)-\mathsf R(c_i)^*$,
\[
\Tr(\rho Z_i)=D^{-1}\|(\mathsf R(c_i)-I)W\|_{HS}^2\ge(\Delta-\theta)^2 .
\]
Let $\mathcal E_i(X)=\frac16\sum_h\mathsf R(h_i)X\mathsf R(h_i)^*$, with $h_i$ the element $h$ placed in factor $i$. These conditional expectations
commute. Since $(I-\Pi)W=0$, $\Tr[(I-\Pi)\mathcal E_i(\rho)]\le\theta^2$. The state $\mathcal E_i(\rho)$ is an average of six states of rank $D$, one
of them $\rho$, so after pinching by $\Pi$ and $I-\Pi$ its two blocks have rank at most $D$ and $5D$. Pinching does not decrease entropy, and
$x\mapsto h_2(x)+x\log5$ increases on $[0,5/6]$, so $S(\mathcal E_i(\rho))\le\log D+h_2(\theta^2)+\theta^2\log5$. For a conditional expectation
$S(\mathcal E(\sigma))-S(\sigma)=D(\sigma\|\mathcal E(\sigma))$. Since the $\mathcal E_j$ commute, monotonicity of relative entropy under
$\mathcal E_{i-1}\cdots\mathcal E_1$ bounds the entropy increment of $\mathcal E_i$, applied after $\mathcal E_{i-1}\cdots\mathcal E_1$, by
$D(\rho\|\mathcal E_i\rho)$. Hence $\bar\rho=\mathcal E_N\cdots\mathcal E_1(\rho)$ satisfies $S(\bar\rho)\le\log D+N[h_2(\theta^2)+\theta^2\log5]$.

The state $\bar\rho$ commutes with $\mathsf R(\Sigma)$, so $\bar\rho=\bigoplus_\lambda p_\lambda(I_{d_\lambda}/d_\lambda\otimes\sigma_\lambda)$ over the
irreducible representations $\lambda$ of $S_3^N$, and $S(\bar\rho)\ge\sum_\lambda p_\lambda\log d_\lambda=\sum_\lambda p_\lambda k_\lambda$, where $k_\lambda$
is the number of factors of $\lambda$ that carry the two-dimensional representation of $S_3$. The operator $Z_i$ is central in the group
algebra, so $\Tr(\bar\rho Z_i)=\Tr(\rho Z_i)$, and it acts as $3I$ on the sectors whose factor $i$ is two-dimensional and as $0$ on the others.
Hence $N(\Delta-\theta)^2\le\sum_i\Tr(\bar\rho Z_i)=3\sum_\lambda p_\lambda k_\lambda\le3S(\bar\rho)$.
\end{proof}

The third fact runs the same count on the three-cycles alone. It needs no relation between the transpositions of different copies. Put
$\chi(x)=h_2(\min\{\frac23x,\frac23\})+\min\{\frac23x,\frac23\}$, a concave function.

\begin{theorem}[Phase-robust sector bound]\label{thm:imp-phase}
Let $c_i,x_i\in U(D)$, $1\le i\le N$, satisfy $c_i^3=x_i^2=I$ and $x_ic_ix_i=c_i^{-1}$, with no condition on how $x_i$ and $x_j$ commute. Put
$\Delta_i=\|c_i-I\|_2$ and $f(g)=c_1^{g_1}\cdots c_N^{g_N}$ for $g\in\Z_3^N$, let $g^{(i)}$ be $g$ with its $i$-th coordinate negated, and, for
uniform $g$, let
\[
\hat a_i^2=\mathbb E_g\|f(g+e_i)-f(g)c_i\|_2^2,\qquad \hat b_i^2=\mathbb E_g\|x_if(g^{(i)})-f(g)x_i\|_2^2 .
\]
Then $\hat a_i^2\le\frac43\sum_{j>i}\|c_ic_j-c_jc_i\|_2^2$, $\hat b_i^2\le\frac43\sum_{j\ne i}\|x_ic_j-c_jx_i\|_2^2$, and
\[
\log D\ \ge\ \sum_{i=1}^N\Big[\frac{(\Delta_i^2-\hat a_i^2)_+}3-\frac{\hat b_i^2}{2\ln2}-\chi(\hat a_i^2)\Big].
\]
\end{theorem}

\begin{proof}
\emph{The ordered-product errors.} For a unitary $h$, $\|(I-\mathcal A_j)h\|_2^2=\frac13\|hc_j-c_jh\|_2^2$ for the average $\mathcal A_j$ of conjugation by
the powers of $c_j$. Moving $c_i$ through the suffix $c_{i+1}^{g_{i+1}}\cdots c_N^{g_N}$ and averaging over its exponents gives
$\hat a_i^2=2-2\,\mathrm{Re}\langle c_i,\mathcal A_N\cdots\mathcal A_{i+1}c_i\rangle$, and Lemma~\ref{lem:projection-product} gives the first bound. For the
second, the maps $\mathcal B_j(Z)=\frac13\sum_kc_j^{-k}Zc_j^{\epsilon_jk}$, with $\epsilon_i=-1$ and $\epsilon_j=1$ for $j\ne i$, are orthogonal
projections. The local relation gives $\mathcal B_i(x_i)=x_i$, and $\|(I-\mathcal B_j)x_i\|_2^2=\frac13\|x_ic_j-c_jx_i\|_2^2$ for $j\ne i$. Moreover
$\hat b_i^2=2-2\,\mathrm{Re}\langle x_i,\mathcal B_N\cdots\mathcal B_1x_i\rangle$, and Lemma~\ref{lem:projection-product} gives the second bound.

\emph{Embedding.} Let $W\xi=3^{-N/2}\sum_gf(g)\xi\otimes|g\rangle$, let $\mathsf C_i$ shift the $i$-th coordinate of $\ell^2(\Z_3^N)$ by $-1$, let
$\mathsf F_i$ negate it, and put $\mathsf X_i=x_i\otimes\mathsf F_i$. The $\mathsf C_i$ commute, $\mathsf X_i\mathsf C_i\mathsf X_i=\mathsf C_i^{-1}$ and
$\mathsf X_i\mathsf C_j\mathsf X_i=\mathsf C_j$ for $j\ne i$. Moreover $D^{-1/2}\|\mathsf C_iW-Wc_i\|_{HS}=\hat a_i$ and
$D^{-1/2}\|\mathsf X_iW-Wx_i\|_{HS}=\hat b_i$.

\emph{Entropy.} Let $\rho=WW^*/D$, and let $\mathcal T_i$ average conjugation by the powers of $\mathsf C_i$. The twirl moves weight at most
$\min\{\frac23\hat a_i^2,\frac23\}$ outside the range of $WW^*$, into a block of rank at most $2D$, so $S(\mathcal T_i\rho)\le\log D+\chi(\hat a_i^2)$. As in
Lemma~\ref{lem:imp-sector} the twirls commute, and the increment of each twirl, the Holevo quantity of the conjugates of the current state, can
only decrease under the earlier twirls by data processing for the Holevo quantity. So $\bar\rho=\mathcal T_1\cdots\mathcal T_N\rho$ has
$S(\bar\rho)\le\log D+\sum_i\chi(\hat a_i^2)$.

\emph{Sectors.} Let $p$ be the joint distribution of the eigenvalue labels $Z_i\in\{1,\omega,\omega^2\}$ of the $\mathsf C_i$, the same for $\rho$
and $\bar\rho$. Then $S(\bar\rho)\ge H(p)\ge\sum_iH(Z_i\mid Z_{-i})$. Let $Q_\alpha$ and $P_\beta$ be the spectral projections of $c_i$ and
$\mathsf C_i$. The numbers $r_{\alpha\beta}=D^{-1}\|P_\beta WQ_\alpha\|_{HS}^2$ form a joint distribution whose marginals are the spectral distribution
of $c_i$ in the normalized trace and that of $\mathsf C_i$ in $\rho$. Distinct cube roots of unity are at squared distance $3$, so
$\hat a_i^2=3\Pr(\alpha\ne\beta)$ and $\Delta_i^2=3\Pr(\alpha\ne1)$, whence $s_i:=\Pr(Z_i\ne1)\ge(\Delta_i^2-\hat a_i^2)_+/3$. Next, the normalized
vectorizations of $W$ and of $\mathsf X_iWx_i^*$ are at squared distance $\hat b_i^2$, so their real overlap is $1-\hat b_i^2/2$. Measuring the joint
spectral projections of the $\mathsf C_j$ gives $p$ from the first and its flip $p^{(i)}$ in coordinate $i$ from the second, because $\mathsf X_i$
exchanges $\omega$ and $\omega^2$ in coordinate $i$ and fixes the other labels. By the Cauchy--Schwarz inequality, outcome by outcome,
$\sum_z\sqrt{p(z)p(z^{(i)})}\ge1-\hat b_i^2/2$. For $0\le x\le1$ one has $1-h_2(x)\le(1-2\sqrt{x(1-x)})/\ln2$: with $y=2x-1$, the difference of
$1-\sqrt{1-y^2}$ and $\ln2\,(1-h_2(x))$ vanishes with its derivative at $y=0$ and has second derivative $(1-y^2)^{-3/2}-(1-y^2)^{-1}\ge0$. So,
for fixed other labels with masses $m,m'$ on the two nontrivial values of $Z_i$, the conditional entropy receives at least
$m+m'-(m+m'-2\sqrt{mm'})/\ln2$. Summing gives $H(Z_i\mid Z_{-i})\ge s_i-\hat b_i^2/(2\ln2)$.
\end{proof}

\subsection{Proof of the counting criterion}

\begin{proof}[Proof of Theorem~\ref{thm:criterion}]
Both parts use the copies $A_i=\phi(w_i)\phi(a)\phi(w_i)^{-1}$ and $B_i=\phi(w_i)\phi(b)\phi(w_i)^{-1}$, $1\le i\le N$, of an assignment $\phi$ extended
to the free group. They have three properties. The $S_3$ relations of each pair are conjugates of $a^2,b^2,(ab)^3$, so they have the defect of a
relator. The product $A_iB_i=\phi(w_i)\phi(ab)\phi(w_i)^{-1}$ has the displacement of $\phi(ab)$. And a commutator of two letters from copies
$i\ne j$, or of their inverses, is a conjugate of $[w_ixw_i^{-1},w_jyw_j^{-1}]^{\pm1}$ with $x,y\in\{a,b\}$, so its area is at most $\Lambda$. A
word of area $A$ is freely a product of $A$ conjugates of relators, so under a unitary assignment it lies within $A$ times the relator defect of
$I$; for permutations the same holds for the normalized Hamming distance, which is bi-invariant.

\emph{Permutations.} Let $f:E\to\Sym_n$ have $f(1)=\mathrm{id}$ and multiply and separate correctly up to $1/r$, with $r\ge10^8N\Lambda(2l-1)$, and put
$\phi(s)=f(s)$ for $s\in S$. For a word $s_1\cdots s_m$ of $R$ with prefixes $p_k$, the chunk contains $p_ks_{k+1}=p_{k+1}$ and $p_m=1$, and
$s\cdot s^{-1}=1$, so $m-1$ merges and at most $m$ inverse replacements give $d(\phi(w),\mathrm{id})\le\delta:=(2l-1)/r$ for every $w\in R$. In the
same way $d(\phi(ab),f(ab))\le1/r$, and $d(f(ab),\mathrm{id})\ge1-1/r$ because $ab\ne1$; so $d(\phi(ab),\mathrm{id})\ge1-2/r\ge\frac12$. Regard
permutations as unitary matrices, with $\|P-I\|_2^2=2d(P,\mathrm{id})$. The $S_3$ relators of each pair then have defect at most $e=\sqrt{2\delta}$,
and Lemma~\ref{lem:imp-s3round} gives exact representations $u_i$ within $10e$ of the normal forms. Commuting two normal forms from different
pairs takes at most nine interchanges of letters, so their commutator has $d\le9\Lambda\delta$ and $\|\cdot\|_2\le3\sqrt{2\Lambda\delta}$, and rounding adds at
most $40e$; so $\eta\le43\sqrt{2\Lambda\delta}$. The element $u_i(st)$ lies within $9e$ of $A_iB_i$, whose displacement in $\|\cdot\|_2$ is at least $1$,
so $\Delta\ge1-9\sqrt{2\delta}\ge\frac12$. Then $\theta^2=2N\eta^2\le7396N\Lambda\delta<10^{-4}$, and Lemma~\ref{lem:imp-sector} gives
$\log n\ge N[(49/100)^2/3-h_2(10^{-4})-10^{-4}\log5]\ge N/16$. So $\log\sigma_E(r)\ge N/16$.

\emph{Unitaries.} Let $\phi$ map every word of $R$ within $e$ of $I$, with $\|\phi(ab)-I\|_2\ge\Delta$ and $K\sqrt N\,\Lambda e\le1$, where
$K=K_\Delta=\sqrt{80000\,(1+\log(1/\Delta))}/\Delta$. Lemma~\ref{lem:imp-s3round} gives exact pairs $x_i=u_i(s)$ and $c_i=u_i(st)$ with
$\|x_i-A_i\|_2\le e$ and $\|c_i-A_iB_i\|_2\le9e$, so $\Delta_i\ge\Delta-9e$. Four and two interchanges of letters give
$\|c_ic_j-c_jc_i\|_2\le4\Lambda e+36e\le40\Lambda e$ and $\|x_ic_j-c_jx_i\|_2\le2\Lambda e+20e\le22\Lambda e$ for $i\ne j$. By Theorem~\ref{thm:imp-phase} and the
concavity of $\chi$,
\[
\log D\ \ge\ N\Big[\frac{(\Delta-9/K)^2-\bar a}{3}-\frac{\bar b}{2\ln2}-\chi(\bar a)\Big],\qquad \bar a\le\frac{3200}{3K^2},\quad \bar b\le\frac{1936}{3K^2},
\]
where $\bar a$ and $\bar b$ are the averages of $\hat a_i^2$ and $\hat b_i^2$, since $e\le1/K$ and $N\Lambda^2e^2\le1/K^2$. It remains to show that the
bracket is at least $\Delta^2/12$. Since $K\ge36/\Delta$, $(\Delta-9/K)^2\ge9\Delta^2/16$, so it suffices that
$\bar a/3+\bar b/(2\ln2)+\chi(\bar a)\le5\Delta^2/48$. Here $\bar a/3+\bar b/(2\ln2)<821.2/K^2$. With $q=2\bar a/3<711.2/K^2$, the bound
$(1-q)\log\frac1{1-q}\le q\log e$ gives $\chi(\bar a)=h_2(q)+q\le q\log(2e/q)$, and $x\log(2e/x)$ increases for $x<2$. Put $L=\log(1/\Delta)$, so that
$K^2=80000(1+L)/\Delta^2$. Multiplying by $K^2$ and using $\log(1+L)\le L+\frac1{10}$ for $L\ge0$, the left side becomes at most
\[
821.2+711.2\Big[\log\frac{2e\cdot80000}{711.2}+\log(1+L)+2L\Big]\ \le\ 7476+2134L\ \le\ \frac5{48}\,80000\,(1+L),
\]
which is the required inequality. Hence $\log D\ge N\Delta^2/12$.
\end{proof}

\subsection{Proof of sequential extraction}

\begin{proof}[Proof of Theorem~\ref{thm:extraction}]
\emph{A test with perfect completeness.} In every round the observer submits half of a fresh maximally entangled pair, measures
$\{\Pi_\Phi,I-\Pi_\Phi\}$ on the output and the reference, keeps the outcome, discards the pair and continues. Let $p_t$ be the probability that
the tests of rounds $1,\dots,t-1$ all pass, $\ell_t$ the probability of failing round $t$ given that, and $F\in\{1,\dots,n,\infty\}$ the first
failure, with $\infty$ for none. The ideal experiment never fails, so $f=\Pr(F\le n)\le\epsilon$. Hence $p_t\ge1-\epsilon$ and
$\sum_tp_t\ell_t=f\le\epsilon$.

\emph{Entropy ledger.} Let $\rho_t$ be the memory state given that the first $t-1$ tests passed; for the bookkeeping, a branch that fails is
frozen. Since the input marginal is $\tau_d$, the round takes $\rho_t$ to $T_t(\rho_t)=\mathrm{Tr}_d\,U_t(\tau_d\otimes\rho_t)U_t^*$, raising the
entropy by $\Delta S_t$. Recording the outcome lowers the average conditional entropy by at most $h_2(\ell_t)$. Exporting $k$ qubits lowers the
entropy by at most $k$, and importing $k$ qubits in a state of entropy $k-c$ raises it by $k-c$, so an exchange lowers it by at most the imported
deficit $c$. The average entropy starts at $Q-P$ and ends at most at $Q$, so
\[
\sum_tp_t\Delta S_t\le P+C+\sum_tp_th_2(\ell_t)\le B+H(F)\le B+h_2(\epsilon)+\epsilon\log n,
\]
where $\sum_tp_th_2(\ell_t)=H(F)$ by the chain rule and $H(F)\le h_2(f)+f\log n$. Therefore
$\sum_tp_t(\ell_t+\Delta S_t)\le B+\epsilon+h_2(\epsilon)+\epsilon\log n$, and since $\sum_tp_t\ge n(1-\epsilon)$, some round has
$\ell_t+\Delta S_t\le t_n$.

\emph{The extracted channel.} Each round's unitary acts only on the input and the $Q$ memory qubits, and the memory state $\rho_t$ conditioned on
the earlier passes does not depend on the current input, since the earlier test registers are discarded. So the round defines a channel $\Psi_t$
with a bath of dimension $2^Q$, and its support leakage is exactly $\ell_t$. For an arbitrary current input with reference, the actual state is
$p_t$ times a normalized pass-block state $\sigma$ plus a failure block, while the ideal state $\omega$ lies in the pass block. Closeness of the
complete experiments gives $\frac12\|p_t\sigma-\omega\|_1+\frac{1-p_t}2\le\epsilon$, hence $p_t\frac12\|\sigma-\omega\|_1\le\epsilon$ and
$\ddia(\Psi_t,\Phi)\le\epsilon/(1-\epsilon)$.
\end{proof}

The device performs no test. The tests and the conditioning on success belong to one particular observer and to the proof.
